\chapter{Parent Company Network and Remote Access Platforms}
\label{app:cato-claroty}
\future{EFC is still adopting these platforms from the parent company. The details below describe the products in general; how they apply to EFC sites depends on how the parent company's IT/OT team has deployed them, and should be confirmed with that team.}

This appendix summarizes the two platforms used by the parent company: \textbf{Cato Networks}, the network and security platform, and \textbf{Claroty xDome Secure Access}, used for remote access to PLCs and other plant equipment. (Claroty is sometimes heard as ``Clarity.'')

\section{Cato Networks (Cato SASE Cloud)}
\subsection{What It Is}
Cato is a Secure Access Service Edge (SASE) platform: a cloud service that replaces traditional site-to-site VPNs, branch firewalls, and remote-access VPNs with a single managed network and security service.\footnote{Cato Networks, Knowledge Base, \emph{Welcome to Cato Networks} and \emph{What is Cato's ZTNA solution}; Cato Networks press release, \emph{Cato Introduces Industry's First SASE-native IoT/OT Security Solution}, December 10, 2024.}
\begin{itemize}[itemsep=1pt]
  \item \textbf{Cato SASE Cloud:} a global network of points of presence (PoPs) that carries traffic between sites and users and applies security inspection: SD-WAN, next-generation firewall, secure web gateway, threat prevention, and zero-trust network access (ZTNA).
  \item \textbf{Cato Socket:} an appliance installed at each site that connects the site to the Cato cloud over encrypted tunnels and provides a local (LAN) firewall.
  \item \textbf{Cato Client:} software for Windows, macOS, Linux, iOS, Android, and Chromebook that connects remote users, with multi-factor authentication and device posture checks. Clientless browser access is also available.
\end{itemize}

\subsection{IoT/OT Features}
In December 2024 Cato added IoT/OT security to the platform:
\begin{itemize}[itemsep=1pt]
  \item automatic discovery and classification of connected devices (type, manufacturer, and version);
  \item policies by device type or group, for example limiting what a PLC or HMI may talk to; and
  \item threat prevention for those devices through DNS security, intrusion prevention, and anti-malware.
\end{itemize}
Cato's published material does not describe industrial-protocol (Modbus, EtherNet/IP) inspection; protocol-level control inside the plant remains the job of the controls network design in Chapter~\ref{ch:networking}.

\subsection{What It Means for Controls}
\begin{itemize}[itemsep=1pt]
  \item The plant and business networks connect to the rest of the company through the Cato Socket. The controls networks in Chapter~\ref{ch:networking} stay behind their own NAT/firewall boundary (or the M580 gateway); Cato does not replace that boundary.
  \item Any connection between the controls network and the plant network or the wider company (SCADA clients, historian replication, remote access) should be agreed with corporate IT as a defined rule, not left open.
  \item Remote engineering access to PLCs should use Claroty (below), which is built for that purpose, rather than general Cato remote access straight into a controls network.
\end{itemize}

\section{Claroty xDome Secure Access}
\subsection{What It Is}
Claroty is an OT cybersecurity company. \textbf{xDome Secure Access} (formerly Claroty SRA) is its remote access product for cyber-physical systems: it lets engineers and vendors reach PLCs, HMIs, and engineering workstations in a plant under control and with a full record.\footnote{Claroty, \emph{xDome Secure Access} product page and solution brief.}
\begin{itemize}[itemsep=1pt]
  \item \textbf{Zero-trust access:} each user is given access to specific devices only, with multi-factor authentication and a credential vault so device passwords are not shared with the user.
  \item \textbf{Approval workflows and just-in-time access:} a user requests a session; an approver (for example the site engineer) grants it, including from a mobile device, for a set time.
  \item \textbf{Session recording and live monitoring:} sessions are logged and can be video-recorded; an administrator can watch a live session and disconnect it.
  \item \textbf{No inbound firewall openings:} the site component connects outward, so no inbound ports are opened into the plant network.
  \item \textbf{Access methods:} browser-based RDP, SSH, VNC, and HTTP/S without installing an agent, plus ``micro-tunnels'' for OT-specific protocols so thick-client engineering software can reach a PLC.
  \item \textbf{Deployment:} as a cloud (SaaS) service or on premises.
\end{itemize}

\subsection{How We Would Use It}
\begin{itemize}[itemsep=1pt]
  \item An EFC engineer or a vendor requests access to a specific PLC or engineering workstation for a stated reason and time window.
  \item After approval, they connect either to an engineering workstation inside the plant (by RDP) that has Control Expert or Productivity Suite installed, or through a micro-tunnel from their own laptop to the PLC.
  \item Online changes made this way still follow Sections~\ref{sec:review-release}, \ref{sec:online-changes}, and~\ref{sec:change-testing}; the session recording becomes part of the change record.
  \item Vendor accounts are time-limited and named, in line with Chapter~9.
\end{itemize}

\section{How the Two Fit Together}
\begin{center}
\begin{tabular}{@{}P{3.68cm}P{6.59cm}P{6.66cm}@{}}
\toprule
\rowcolor{tblhead} & \textbf{Cato Networks} & \textbf{Claroty xDome Secure Access} \\
\midrule
Layer & Company network and security (WAN, firewall, user access) & OT remote access to specific plant devices \\
Main users & Corporate IT; all company users & Controls engineers, maintenance, and vendors \\
At the site & Cato Socket at the plant network edge & Site component inside the plant, connecting outward \\
Controls role & Carries and protects traffic to and from the plant & Controls, approves, and records sessions to PLCs and HMIs \\
\bottomrule
\end{tabular}
\end{center}
See also Sections~\ref{sec:remote-access} (Remote Access) and~\ref{sec:hardening} (System Hardening) and Chapter~9 (Accounts and Access Control).
